\section{Evaluasi}
\label{sec:evaluasi-bab6}

\subsection{Soal Pilihan Ganda}

\begin{enumerate}
    \item Rumus yang benar untuk menghitung alamat fisik dari
          pasangan \eng{segment:offset} adalah \dots
          \begin{enumerate}
              \item $\text{segmen} + \text{offset}$
              \item $(\text{segmen} \times \hex{10}) +
                    \text{offset}$
              \item $(\text{segmen} \times \dec{20}) +
                    \text{offset}$
              \item $(\text{offset} \times \hex{10}) +
                    \text{segmen}$
          \end{enumerate}

    \item Nilai segmen \hex{1234} menunjuk alamat dasar \dots
          \begin{enumerate}
              \item \hex{1234} \item \hex{12340}
              \item \hex{01234} \item \hex{23400}
          \end{enumerate}

    \item Alamat fisik yang ditunjuk \code{B800h:0000h} adalah
          \dots
          \begin{enumerate}
              \item \hex{B800} \item \hex{0B800}
              \item \hex{B8000} \item \hex{BB000}
          \end{enumerate}

    \item Satu paragraf pada skema pengalamatan 8086 berukuran
          \dots
          \begin{enumerate}
              \item 16 byte \item 256 byte
              \item 1 KB \item 64 KB
          \end{enumerate}

    \item Jumlah pasangan \eng{segment:offset} terbanyak yang dapat
          menunjuk satu alamat fisik adalah \dots
          \begin{enumerate}
              \item 16 \item 64
              \item 1.024 \item 4.096
          \end{enumerate}
\end{enumerate}

\subsection{Soal Uraian}

\begin{enumerate}
    \item Jelaskan mengapa segmentasi diperlukan pada 8086. Kaitkan
          jawaban dengan lebar register segmen, lebar bus alamat,
          dan besar memori yang ingin dijangkau.

    \item Hitung alamat fisik dari pasangan berikut dan tuliskan
          setiap langkahnya.
          \begin{enumerate}
              \item \code{0000h:00A0h}
              \item \code{2000h:0010h}
              \item \code{9FFFh:0010h}
              \item \code{B800h:0080h}
          \end{enumerate}

    \item Jelaskan pengertian \eng{overlapping segmen}. Hitung
          banyaknya pasangan yang menunjuk alamat fisik
          \hex{00040}, tuliskan seluruh pasangan tersebut, dan
          jelaskan mengapa jumlahnya tidak pernah melebihi 4.096.

    \item Gambarkan pembagian memori 1 MB dan jelaskan kegunaan
          empat area: memori konvensional, buffer video, ROM
          adaptor, dan ROM BIOS. Jelaskan pula mengapa alamat
          \hex{B8000} termasuk area video.
\end{enumerate}

\subsection{Kisi-kisi Penilaian}

\begin{table}[H]
    \centering
    \caption{Kisi-kisi evaluasi Bab~\ref{chap:segmentasi}.}
    \label{tab:kisi-bab6}
    \small
    \begin{tabular}{@{}p{1.5cm}p{5.0cm}p{3.0cm}c@{}}
        \toprule
        \textbf{Sub-CPMK} & \textbf{Indikator yang diukur} &
        \textbf{Bentuk soal} & \textbf{Bobot} \\
        \midrule
        1.4 & Menjelaskan alasan segmentasi & Uraian 1 & 20\% \\
        \addlinespace
        1.4 & Menghitung alamat fisik &
        Pilihan ganda 1--3; Uraian 2 & 35\% \\
        \addlinespace
        1.4 & Menjelaskan paragraf dan \eng{overlapping} &
        Pilihan ganda 4, 5; Uraian 3 & 25\% \\
        \addlinespace
        1.4 & Menjelaskan pembagian memori 1 MB & Uraian 4 & 20\% \\
        \bottomrule
    \end{tabular}
\end{table}

\begin{checklistbox}
\vspace{1.1em}
\textbf{Checklist Pencapaian Kompetensi Bab~\ref{chap:segmentasi}}

Seorang mahasiswa dinyatakan kompeten pada Sub-CPMK 1.4 apabila:
\begin{itemize}
    \item[\checked] dapat menjelaskan mengapa 8086 memakai pasangan
          \eng{segment:offset};
    \item[\checked] dapat menghitung alamat fisik dengan menuliskan
          langkahnya;
    \item[\checked] dapat menjelaskan paragraf 16 byte dan segmen
          64 KB;
    \item[\checked] dapat menghitung banyaknya pasangan yang
          menunjuk satu alamat fisik;
    \item[\checked] dapat menyebutkan pembagian memori 1 MB;
    \item[\checked] dapat memverifikasi alamat dengan DEBUG atau
          emu8086.
\end{itemize}
\end{checklistbox}
